\documentclass[aspectratio=169, 12pt]{beamer}

% --- Paquetes (compatible con pdfLaTeX) ---
\usepackage[utf8]{inputenc}
\usepackage[T1]{fontenc}
\usepackage[spanish]{babel}
\usepackage{amsmath}
\usepackage{booktabs}
\usepackage{tikz}
\usetikzlibrary{calc, positioning, shadows.blur}

% --- Colores ---
\definecolor{bg}{HTML}{0F172A}
\definecolor{bg2}{HTML}{1E293B}
\definecolor{red1}{HTML}{EF4444}
\definecolor{yellow1}{HTML}{F59E0B}
\definecolor{green1}{HTML}{22C55E}
\definecolor{txt}{HTML}{F8FAFC}
\definecolor{txt2}{HTML}{94A3B8}
\definecolor{txt3}{HTML}{64748B}
\definecolor{border}{HTML}{334155}

% --- Tema minimalista oscuro ---
\usetheme{default}
\usecolortheme{default}
\setbeamercolor{background canvas}{bg=bg}
\setbeamercolor{normal text}{fg=txt}
\setbeamercolor{frametitle}{fg=txt}
\setbeamercolor{item}{fg=yellow1}
\setbeamercolor{subitem}{fg=red1}
\setbeamercolor{block title}{fg=yellow1, bg=bg2}
\setbeamercolor{block body}{fg=txt2, bg=bg2}
\setbeamercolor{block title alerted}{fg=red1, bg=bg2}
\setbeamercolor{block body alerted}{fg=txt2, bg=bg2}

% --- Sin navegación, footer minimalista ---
\setbeamertemplate{navigation symbols}{}
\setbeamertemplate{frametitle}{
  \vspace{8pt}
  {\large\bfseries\insertframetitle}
}
\setbeamertemplate{footline}{
  \hbox to \paperwidth{
    \hspace{14pt}
    {\tiny\color{txt3}\textbf{TIENDITACAMPUS}}
    \hfill
    {\tiny\color{txt3}\insertframenumber\,/\,\inserttotalframenumber}
    \hspace{14pt}
  }
  \vspace{4pt}
}

% --- Items limpios ---
\setbeamertemplate{itemize items}[circle]
\setbeamertemplate{itemize subitem}[triangle]
\setlength{\leftmargini}{12pt}

% --- Helpers ---
\newcommand{\tag}[1]{%
  \tikz[baseline=(t.base)]{\node[fill=bg2, draw=yellow1!30, text=yellow1,
  rounded corners=8pt, inner xsep=6pt, inner ysep=2pt,
  font=\tiny\bfseries] (t) {\MakeUppercase{#1}};}%
}
\newcommand{\pill}[1]{%
  \tikz[baseline=(p.base)]{\node[fill=red1!15, draw=red1!30, text=red1!80!white,
  rounded corners=12pt, inner xsep=8pt, inner ysep=3pt,
  font=\scriptsize\bfseries] (p) {#1};}%
}
\newcommand{\badge}[2]{%
  \tikz[baseline=(b.base)]{\node[fill=#1!15, text=#1!80!white,
  rounded corners=4pt, inner xsep=4pt, inner ysep=2pt,
  font=\tiny\bfseries] (b) {#2};}%
}

\begin{document}

% ===== SLIDE 1: PORTADA =====
{
\setbeamertemplate{footline}{}
\begin{frame}[plain]
\centering\vspace{0.8cm}

{\tiny\color{txt3}\MakeUppercase{Universidad Polit\'ecnica de Chiapas}}\\[2pt]
{\scriptsize\color{txt2}Proyecto Integrador 2 \,$\cdot$\, Corte 1 \,$\cdot$\, Unidad 2}\\[18pt]

{\fontsize{36}{40}\selectfont\bfseries\color{red1}TienditaCampus}\\[6pt]
{\small\color{txt2}Herramientas digitales para vendedores universitarios}\\[18pt]

\pill{Marco Te\'orico}\hspace{6pt}
\pill{Marco Legal y Normativo}\hspace{6pt}
\pill{Poblaci\'on y Muestra}\\[18pt]

{\footnotesize\color{txt2}
\textbf{H\'ector Isaac Espinoza Mendoza}\\[1pt]
\textbf{Luis Emilio Jaras S\'anchez}\\[1pt]
\textbf{Jesel Moreno Z\'u\~niga}}\\[10pt]

{\tiny\color{txt3}Docente: Roc\'io Crystal Hern\'andez Camacho \,$\cdot$\, Marzo 2026}
\end{frame}
}


\begin{frame}{\tag{Agenda} \; Contenido de la {\color{red1}Presentaci\'on}}
\vspace{4pt}
\begin{columns}[T]
\begin{column}{0.48\textwidth}
\begin{block}{1. Marco Te\'orico}
\begin{itemize}\small
  \item Comercio electr\'onico universitario
  \item Aplicaciones similares
  \item PWA y stack tecnol\'ogico
  \item Arquitectura 3-Tier en AWS
  \item Inventarios perecederos
\end{itemize}
\end{block}
\end{column}
\begin{column}{0.48\textwidth}
\begin{block}{2. Marco Legal y Normativo}
\begin{itemize}\small
  \item LFPDPPP (Datos Personales)
  \item Protecci\'on al Consumidor
  \item NOM-151 y NOM-251
\end{itemize}
\end{block}
\vspace{4pt}
\begin{alertblock}{3. Poblaci\'on y Muestra}
\begin{itemize}\small
  \item Caso pr\'actico integrador
  \item Criterios de selecci\'on
  \item C\'alculo de muestra
\end{itemize}
\end{alertblock}
\end{column}
\end{columns}
\end{frame}


\begin{frame}{\tag{Marco Te\'orico} \; Comercio Electr\'onico {\color{red1}Universitario}}
\vspace{4pt}
\begin{columns}[T]
\begin{column}{0.48\textwidth}
\begin{block}{Definici\'on}
\small El e-commerce es la compra y venta de bienes a trav\'es de Internet (Laudon \& Traver, 2022). En universidades, surge como respuesta al emprendimiento estudiantil con herramientas accesibles.
\end{block}
\begin{block}{Contexto}
\small Seg\'un ANUIES (2024), el \textbf{38\%} de los universitarios mexicanos han vendido informalmente en su campus, pero carecen de herramientas digitales para gestionar sus micro-negocios.
\end{block}
\end{column}
\begin{column}{0.48\textwidth}
\begin{alertblock}{TienditaCampus}
\small PWA que digitaliza la gesti\'on de vendedores universitarios:
\begin{itemize}\small
  \item Control de inventario
  \item Seguimiento de ventas diarias
  \item An\'alisis de rentabilidad
  \item Reducci\'on de p\'erdidas (M\'etodo IQR)
  \item Desde el celular, sin instalar app
\end{itemize}
\end{alertblock}
\end{column}
\end{columns}
\end{frame}


\begin{frame}{\tag{Marco Te\'orico} \; Aplicaciones {\color{red1}Similares} y Nuestra Diferencia}
\vspace{2pt}
\begin{columns}[T]
\begin{column}{0.48\textwidth}
\begin{block}{Square y Toast}
\footnotesize Terminales POS con comisiones del \textbf{3\%+} por transacci\'on. Un estudiante vendiendo tortas de \$25 no puede absorber ese costo.
\end{block}
\vspace{2pt}
\begin{block}{F\'aciLUNCH e Infood}
\footnotesize Apps de pedidos universitarios. Distribuyen comida pero \textbf{no ayudan} a saber cu\'anto producto comprar.
\end{block}
\vspace{2pt}
\begin{block}{Too Good To Go}
\footnotesize Vende sobras a precio rebajado. Soluci\'on \textbf{reactiva}: el vendedor ya perdi\'o dinero.
\end{block}
\end{column}
\begin{column}{0.48\textwidth}
\begin{alertblock}{TienditaCampus}
\textbf{Enfoque PREVENTIVO con M\'etodo IQR:}
\begin{itemize}\small
  \item Analiza \textbf{ventas pasadas}
  \item Filtra \textbf{picos irreales} (outliers)
  \item \textbf{Sugiere cu\'anto comprar}
  \item \textbf{Gratuito}, sin comisiones
  \item Accesible desde el celular
\end{itemize}
\vspace{4pt}
{\color{green1}\bfseries\small Ataca el problema ANTES de que suceda}
\end{alertblock}
\end{column}
\end{columns}
\end{frame}


\begin{frame}{\tag{Marco Te\'orico} \; Aplicaciones Web {\color{red1}Progresivas (PWA)}}
\vspace{4pt}
\begin{columns}[T]
\begin{column}{0.48\textwidth}
\begin{block}{\textquestiondown Qu\'e es una PWA?}
\small Combina lo mejor de apps web y nativas (Google, 2023). Instalables, funcionan offline, se actualizan autom\'aticamente con Service Workers.
\end{block}
\end{column}
\begin{column}{0.48\textwidth}
\begin{block}{Ventajas}
\begin{itemize}\small
  \item \textbf{Sin App Store:} solo un link
  \item \textbf{Offline:} registrar ventas sin red
  \item \textbf{Liviana:} no ocupa espacio
  \item \textbf{Multi-plataforma}
\end{itemize}
\end{block}
\end{column}
\end{columns}
\vspace{6pt}
\centering
\badge{yellow1}{Next.js 14}\hspace{2pt}
\badge{red1}{NestJS}\hspace{2pt}
\badge{blue}{PostgreSQL 16}\hspace{2pt}
\badge{green1}{MongoDB 7}\hspace{2pt}
\badge{violet}{Docker}\hspace{2pt}
\badge{pink}{Nginx}\hspace{2pt}
\badge{teal}{Tailwind CSS v4}\hspace{2pt}
\badge{orange}{AWS EC2}
\end{frame}


\begin{frame}{\tag{Marco Te\'orico} \; Arquitectura {\color{red1}3 Niveles} en AWS}
\vspace{4pt}
\footnotesize Cada capa corre en una instancia EC2 independiente (Fowler, 2002).\\[6pt]
\begin{columns}[T]
\begin{column}{0.32\textwidth}
\begin{block}{Presentaci\'on}
\scriptsize \textbf{Next.js + Nginx}
\begin{itemize}\scriptsize
  \item Interfaz PWA
  \item Reverse proxy
  \item SSR
\end{itemize}
\end{block}
\end{column}
\begin{column}{0.32\textwidth}
\begin{block}{L\'ogica}
\scriptsize \textbf{NestJS API REST}
\begin{itemize}\scriptsize
  \item JWT + OAuth
  \item Validaci\'on DTOs
  \item IP privada
\end{itemize}
\end{block}
\end{column}
\begin{column}{0.32\textwidth}
\begin{alertblock}{Datos}
\scriptsize \textbf{PostgreSQL 16} (relacional) + \textbf{MongoDB 7} (logs). Arquitectura poliglota.
\end{alertblock}
\end{column}
\end{columns}
\end{frame}


\begin{frame}{\tag{Marco Te\'orico} \; Inventarios {\color{red1}Perecederos}}
\vspace{4pt}
\begin{columns}[T]
\begin{column}{0.48\textwidth}
\begin{block}{Problem\'atica}
\small Seg\'un la FAO (2023), hasta el \textbf{30\%} del inventario en micronegocios de alimentos se pierde por falta de control. Los vendedores universitarios son especialmente vulnerables.
\end{block}
\end{column}
\begin{column}{0.48\textwidth}
\begin{block}{Soluci\'on}
\begin{itemize}\small
  \item \textbf{Registro diario} de entradas/salidas
  \item \textbf{Alertas} de caducidad
  \item \textbf{Rentabilidad:} costo vs.\ ganancia
  \item \textbf{Mermas:} patrones de p\'erdida
\end{itemize}
\end{block}
\end{column}
\end{columns}
\vspace{6pt}
\begin{block}{Modelo de Datos}
\small \textbf{Products} $\rightarrow$ \textbf{InventoryRecords} $\rightarrow$ \textbf{DailySales} $\rightarrow$ \textbf{SaleDetails} $\rightarrow$ \textbf{Orders}. Trazabilidad completa, m\'argenes reales (Chopra \& Meindl, 2020).
\end{block}
\end{frame}


\begin{frame}{\tag{Marco Legal} \; Protecci\'on de {\color{red1}Datos Personales}}
\vspace{4pt}
\begin{columns}[T]
\begin{column}{0.48\textwidth}
\begin{alertblock}{LFPDPPP (DOF, 2010)}
\begin{itemize}\small
  \item \textbf{Art.\ 6:} Tratamiento l\'icito e informado
  \item \textbf{Art.\ 9:} Consentimiento expreso para datos sensibles
  \item \textbf{Art.\ 16:} Aviso de privacidad obligatorio
\end{itemize}
\end{alertblock}
\end{column}
\begin{column}{0.48\textwidth}
\begin{block}{Cumplimiento}
\begin{itemize}\small
  \item \textbf{Argon2} para hashing
  \item \textbf{JWT} sin datos sensibles en servidor
  \item \textbf{Google OAuth} delegado
  \item \textbf{.env} cero credenciales en c\'odigo
  \item \textbf{Roles} con m\'inimo privilegio
\end{itemize}
\end{block}
\end{column}
\end{columns}
\end{frame}


\begin{frame}{\tag{Marco Legal} \; Comercio Digital y {\color{red1}Normativas}}
\vspace{4pt}
\begin{columns}[T]
\begin{column}{0.48\textwidth}
\begin{block}{Protecci\'on al Consumidor}
\small \textbf{Cap.\ VIII Bis} --- Transacciones electr\'onicas.
\begin{itemize}\small
  \item \textbf{Art.\ 76 Bis:} Identidad del proveedor
  \item \textbf{Frac.\ V:} Protecci\'on de datos
\end{itemize}
\end{block}
\vspace{2pt}
\begin{block}{NOM-151-SCFI-2002}
\small Conservaci\'on de mensajes de datos. MongoDB como sistema de auditor\'ia y trazabilidad.
\end{block}
\end{column}
\begin{column}{0.48\textwidth}
\begin{block}{NOM-251-SSA1-2009}
\small Higiene alimentaria. El sistema obliga orden \textbf{FIFO}: lo primero que entra es lo primero que se vende.
\end{block}
\vspace{2pt}
\begin{alertblock}{UPChiapas}
\small Alineado con las pol\'iticas de emprendimiento e innovaci\'on, dentro del \textbf{Modelo Educativo Basado en Competencias (MEBC)}.
\end{alertblock}
\end{column}
\end{columns}
\end{frame}


\begin{frame}{\tag{Poblaci\'on y Muestra} \; Caso Pr\'actico {\color{red1}Integrador}}
\vspace{6pt}
\centering
\begin{tikzpicture}
  \node[fill=bg2, draw=border, rounded corners=6pt, minimum width=12.5cm, 
        minimum height=0.9cm, text width=12cm, align=left] at (0, 0) {
    {\color{yellow1}\bfseries\scriptsize TEMA}\hspace{8pt}
    {\color{txt}\footnotesize Gesti\'on del comercio informal de snacks universitarios mediante una app web.}
  };
  \node[fill=bg2, draw=border, rounded corners=6pt, minimum width=12.5cm, 
        minimum height=0.9cm, text width=12cm, align=left] at (0,-1.2) {
    {\color{yellow1}\bfseries\scriptsize DISE\~NO}\hspace{8pt}
    {\color{txt}\footnotesize Cuasi-experimental, mediciones pre-post intervenci\'on, enfoque cuantitativo.}
  };
  \node[fill=bg2, draw=border, rounded corners=6pt, minimum width=12.5cm, 
        minimum height=0.9cm, text width=12cm, align=left] at (0,-2.4) {
    {\color{yellow1}\bfseries\scriptsize POBLACI\'ON}\hspace{8pt}
    {\color{txt}\footnotesize Todos los vendedores estudiantiles de snacks en la UPChiapas.}
  };
  \node[fill=red1!10, draw=red1!30, rounded corners=6pt, minimum width=12.5cm, 
        minimum height=1.1cm, text width=12cm, align=left] at (0,-3.7) {
    {\color{red1}\bfseries\scriptsize MUESTRA}\hspace{8pt}
    {\color{txt}\footnotesize Muestreo intencional: vender productos preparados perecederos al menos \textbf{3 d\'ias/semana}.}
  };
\end{tikzpicture}
\end{frame}


\begin{frame}{\tag{Poblaci\'on y Muestra} \; Criterios de {\color{red1}Selecci\'on}}
\vspace{4pt}
\begin{columns}[T]
\begin{column}{0.48\textwidth}
\begin{block}{Inclusi\'on}
\begin{itemize}\small
  \item Estudiante \textbf{activo} en UPChiapas
  \item Vendedor de \textbf{snacks} en campus
  \item M\'inimo \textbf{1 mes} de experiencia
  \item Vender \textbf{3+ d\'ias/semana}
  \item Disponibilidad para usar la plataforma
  \item Acceso a dispositivo con navegador
\end{itemize}
\end{block}
\end{column}
\begin{column}{0.48\textwidth}
\begin{alertblock}{Exclusi\'on}
\begin{itemize}\small
  \item Vendedores de \textbf{comidas completas}
  \item \textbf{Ambulantes} sin punto fijo
  \item \textbf{Ajenos} a la comunidad universitaria
  \item Personal administrativo o docente
  \item Negocios con sistemas contables
\end{itemize}
\end{alertblock}
\end{column}
\end{columns}
\end{frame}


\begin{frame}{\tag{Poblaci\'on y Muestra} \; C\'alculo de la {\color{red1}Muestra}}
\vspace{4pt}
\begin{block}{}
\centering
\large $n = \dfrac{N \cdot Z^2 \cdot p \cdot q}{e^2 (N-1) + Z^2 \cdot p \cdot q}$
\end{block}

\vspace{4pt}
\begin{columns}[c]
\begin{column}{0.22\textwidth}\centering
{\Large\bfseries\color{red1}N=150}\\{\tiny\color{txt2}Poblaci\'on}
\end{column}
\begin{column}{0.22\textwidth}\centering
{\Large\bfseries\color{red1}Z=1.96}\\{\tiny\color{txt2}Confianza 95\%}
\end{column}
\begin{column}{0.22\textwidth}\centering
{\Large\bfseries\color{red1}p=0.5}\\{\tiny\color{txt2}Proporci\'on}
\end{column}
\begin{column}{0.22\textwidth}\centering
{\Large\bfseries\color{red1}e=5\%}\\{\tiny\color{txt2}Error}
\end{column}
\end{columns}

\vspace{8pt}
\centering
\tikz{\node[fill=red1!15, draw=red1!30, rounded corners=8pt, inner sep=8pt]{
  {\Huge\bfseries\color{red1}n $\approx$ 108}\quad
  {\small\color{txt2}vendedores a encuestar}
};}

\vspace{6pt}
\small\color{txt2}\textbf{Muestreo intencional} $+$ bola de nieve durante recesos y cambios de clase.
\end{frame}


{
\setbeamertemplate{footline}{}
\begin{frame}[plain]
\centering\vspace{1.5cm}

{\fontsize{36}{40}\selectfont\bfseries\color{red1}\textexclamdown Gracias!}\\[10pt]
{\color{txt2}\textquestiondown Preguntas o comentarios?}\\[20pt]

\pill{tienditacampus.com}\\[18pt]

{\footnotesize\color{txt2}
\textbf{H\'ector Isaac Espinoza Mendoza}\\[1pt]
\textbf{Luis Emilio Jaras S\'anchez}\\[1pt]
\textbf{Jesel Moreno Z\'u\~niga}}\\[12pt]

{\tiny\color{txt3}Universidad Polit\'ecnica de Chiapas \,$\cdot$\, Proyecto Integrador 2 \,$\cdot$\, Marzo 2026}
\end{frame}
}

\end{document}
